\section{\MakeTextUppercase{Заголовок второй главы}}

\subsection{Заголовок раздела}

Вводный текст раздела.

\subsubsection{Заголовок подраздела}

Текст первого подраздела.

\begin{definition}\label{def:yalv}
	Формулой языка логики высказываний называется \ldots
\end{definition}

Опираясь на определение~\ref{def:yalv}, сформулируем и докажем теорему~\ref{thm:tautology_criteria}.

\begin{theorem}\label{thm:tautology_criteria}
	Формула~ЯЛВ является тавтологией тогда и только тогда, когда \ldots
\end{theorem}

Для этого нам понадобится лемма~\ref{lemma:coliear_vectors}.

\begin{lemma}\label{lemma:coliear_vectors}
	 Если вектора $\vec{a}$~и~$\vec{b}$ коллинеарны, то \ldots
\end{lemma}

\subsubsection{Заголовок следующего подраздела}

Текст второго подраздела.

Примеры круговых диаграмм, построенных в~TeX, представлены на рисунках~\ref{fig:start_ank_1}~и~\ref{fig:start_ank_2}.

\begin{figure}[!h]
	\minipage{0.48\textwidth}\centering
	\textbf{\small 1. Интересно ли Вам изучать математику?}\vspace*{-1em}
	
	\begin{tikzpicture}[
		pie chart,
		slice type={ans1}{yellow!90!black},
		slice type={ans2}{green!60!black},
		slice type={ans3}{blue!60},
		slice type={ans4}{red!70},
		slice type={ans5}{gray!70},
		pie values/.style={font={\footnotesize}},
		scale=2.25
		]
		\pie{}{3.4/ans4,17.2/ans3,48.3/ans2,31.0/ans1}
		
		\legend[shift={(-47pt,-30pt)}]{{\footnotesize Да, интересно}/ans1, {\footnotesize Скорее да, чем нет}/ans2, {\footnotesize Скорее нет, чем да}/ans3}
		\legend[shift={(5pt,-30pt)}]{{\footnotesize Нет, не интересно}/ans4, {\footnotesize Затрудняюсь ответить}/ans5}
	\end{tikzpicture}\vspace*{27pt}
	\caption{Ответы на вопрос №~1}\label{fig:start_ank_1}
	\endminipage\hfill
	\minipage{0.48\textwidth}\centering
	\textbf{\small 2. Какой предмет Вам больше нравится: алгебра или геометрия?}\vspace*{-1em}
	
	\begin{tikzpicture}[
		pie chart,
		slice type={ans1}{yellow!90!black},
		slice type={ans2}{green!60!black},
		slice type={ans3}{blue!60},
		slice type={ans4}{red!70},
		slice type={ans5}{gray!70},
		pie values/.style={font={\footnotesize}},
		scale=2.25
		]
		\pie{}{10.3/ans5,10.3/ans4,13.8/ans3,13.8/ans2,51.7/ans1}
		
		\legend[shift={(-42pt,-30pt)}]{{\footnotesize Алгебру люблю больше, чем геометрию}/ans1, {\footnotesize Геометрию люблю больше, чем алгебру}/ans2, {\footnotesize Одинаково люблю и алгебру, и геометрию}/ans3, {\footnotesize Не очень люблю оба этих предмета}/ans4, {\footnotesize Затрудняюсь ответить}/ans5}
	\end{tikzpicture}
	\caption{Ответы на вопрос №~2}\label{fig:start_ank_2}
	\endminipage
\end{figure}

Также приведём пример столбчатой диаграммы.

\begin{figure}[!h]\centering
	\textbf{\small 6. С какими из перечисленных видов теорем Вы сталкивались на уроках математики?}\vspace*{10pt}
	
	\begin{tikzpicture}[scale=6]\footnotesize
		\draw[thin, color=gray] (0,0)--(2,0);
		\draw[thin, color=gray] (0,0.1)--(2,0.1);
		\draw[thin, color=gray] (0,0.2)--(2,0.2);
		\draw[thin, color=gray] (0,0.3)--(2,0.3);
		\draw[thin, color=gray] (0,0.4)--(2,0.4);
		\draw[thin, color=gray] (0,0.5)--(2,0.5);
		\draw[thin, color=gray] (0,0.6)--(2,0.6);
		\draw[thin, color=gray] (0,0.7)--(2,0.7);
		\draw[thin, color=gray] (0,0.8)--(2,0.8);
		\draw[thin, color=gray] (0,0.9)--(2,0.9);
		\draw[thin, color=gray] (0,1)--(2,1);
		\fill[blue!50] (0.15,0) rectangle (0.35,0.897);
		\fill[blue!50] (0.65,0) rectangle (0.85,0.897);
		\fill[blue!50] (1.15,0) rectangle (1.35,0.828);
		\fill[blue!50] (1.65,0) rectangle (1.85,0.586);
		\node[anchor=east, color=gray] at (0,0) {0\,\%};
		\node[anchor=east, color=gray] at (0,0.1) {10\,\%};
		\node[anchor=east, color=gray] at (0,0.2) {20\,\%};
		\node[anchor=east, color=gray] at (0,0.3) {30\,\%};
		\node[anchor=east, color=gray] at (0,0.4) {40\,\%};
		\node[anchor=east, color=gray] at (0,0.5) {50\,\%};
		\node[anchor=east, color=gray] at (0,0.6) {60\,\%};
		\node[anchor=east, color=gray] at (0,0.7) {70\,\%};
		\node[anchor=east, color=gray] at (0,0.8) {80\,\%};
		\node[anchor=east, color=gray] at (0,0.9) {90\,\%};
		\node[anchor=east, color=gray] at (0,1) {100\,\%};
		\node[anchor=north] at (0.25,0) {Теоремы-свойства};
		\node[anchor=north] at (0.75,0) {Теоремы-признаки};
		\node[anchor=north, align=center] at (1.25,0) {Теоремы\\ существования};
		\node[anchor=north, align=center] at (1.75,0) {Теоремы вида\\ <<не более чем один>>};
		\node[anchor=south] at (0.25,0.897) {89,7\,\%};
		\node[anchor=south] at (0.75,0.897) {89,7\,\%};
		\node[anchor=south] at (1.25,0.828) {82,8\,\%};
		\node[anchor=south] at (1.75,0.586) {58,6\,\%};
	\end{tikzpicture}
	\caption{Ответы на вопрос №~6}\label{fig:start_ank_6}
\end{figure}

\newpage

На этом подраздел завершён.

\newpage

\subsection{Заголовок следующего раздела}

Наконец, приведём пример листинга кода:

\lstinputlisting[language=C++]{listings/helloworld.cpp}

Листинги можно нумеровать, называть и ссылаться на них. Например, рассмотрим листинг~\ref{list:hellouser}.

\lstinputlisting[language=C++, numbers=none, caption={Пример программы, осуществляющей ввод-вывод}, label={list:hellouser}]{listings/hellouser.cpp}